\section{Sparse pivot scheduling for primitive cocycle discovery}
\label{sec:cocycle-sparse}

Section~\ref{sec:cocycle-seeds} separates cocycle discovery from span
optimization. The former is another significant cost: its initial solver
processes triangular-face equations in tetrahedron order and always pivots
on the leftmost remaining column. Even sparse input can then create dense
intermediate rows and nonunit rational pivots. We now retain the same
maximal-tree gauge and all face equations, but change their elimination order.

\subsection{Exact elimination with live column incidence}

Maintain every active row as a sparse dictionary and every column as the
set of active rows containing it. Buckets indexed by row support size select
a shortest nonzero row; row index breaks ties. Within that row, prefer a
coefficient of absolute value one, then a column appearing in fewer active
rows, then the column index. These choices are deterministic. They aim to
reduce fill and rational arithmetic; they are not a theorem that either
quantity remains small on all knot exteriors.

Normalize the chosen row at its pivot $p$ and store it as an equation
\[
 x_p=-\sum_{i\ne p}a_i x_i.
\]
Remove that row from all active column sets. Only rows in the incidence set
of $p$ need an elimination update. For each one, subtract its pivot
coefficient times the normalized row, maintaining exact rational values,
column incidence and support-size buckets. A row that becomes zero is
retired. Repeated global edges in a face were already combined before
this procedure, so no simplicial independence is assumed.

The invariant is that stored pivot equations and active rows have exactly
the original rational nullspace. Normalization and row subtraction are
invertible row operations. Every active row has zero entries in previously
eliminated columns. Thus a stored pivot equation only refers to variables
eliminated later or left free. Reverse \emph{elimination order}, rather than
reverse column order, gives valid back substitution. The number of free
variables is the exact nullity; the API still requires one.

Clearing denominators and dividing by the coordinate gcd gives the primitive
integer generator as before. Require its last nonzero chord coordinate to
be positive. This recovers the initial solver's sign convention: in its
leftmost-pivot echelon form, the only free coordinate is the last nonzero
coordinate of the kernel vector. Since the tree gauge, chord ordering and
primitive sign are unchanged, the new method returns the identical heights
and normal coordinates, not merely a homologous representative. Consequently
the downstream span problem and source-disc certificate remain the same
whenever the shared budget lets both versions finish.

There are $O(T)$ rows and columns. Active rows, column incidence and stored
pivot equations can require $O(T^2)$ entries. A conservative dense bound is
still $O(T^3)$ exact arithmetic updates, with polynomial bit cost controlled
by minors of the original sparse integer matrix. Shortest-bucket scans,
deterministic row selection and incidence ordering do not improve that
worst-case bound. The implementation exposes peak active-row support and
nonunit pivot counts to distinguish observed sparsity from a global theorem.
Cooperative work guards cover initialization, pivot normalization, incidence
updates and back substitution; they remain an operational allowance rather
than a bit-complexity measure.

\subsection{Preserved source witnesses and observed arithmetic savings}

An independent dense rational elimination test checks 240 small integer
matrices with varied support, rank and coefficients, including a binary
scaling case. It compares nullity decisions and rank-one kernels against
full reduced row echelon form. Existing geometry tests include relabeled
layered solid tori, a higher-rank rejection, and a rank-one manifold with
homological torsion that still requires exact nonunit pivots. They also
retain source-proof mutations, cancellation, shared-budget endpoints and
the genus-one unknot miss.
All sixteen focused tests pass in 3.372 seconds; all 1,166 maintained tests
pass in 237.520 seconds. The source audit runs concurrently with the latter
and takes 60.221 seconds; neither elapsed time is a performance comparison.

The pinned source audit compares the published solver at commit
\code{7a13523cc} with the new implementation on all 84 earlier diagrams.
Every vertex label, local height and normal coordinate is identical.
All seventeen positive source queries replay independently, and their
certificate hashes match the prior stored proofs. Sixty-two queries finish
without a disc in the tested family, and the same five reach the shared
two-million-guard limit. This scheduling change does not add source coverage.
The previous Regina comparisons apply to these identical vectors; they are
not counted as new external runs.

Across the 84 preparations, elimination updates fall from 2,286,028 to
286,748 and nonunit pivots from 861 to zero. The largest active support is
ten. For Gordian, updates fall from 1,124,538 to 72,167 and preprocessing
guards from 1,564,758 to 208,461. Nevertheless its subsequent raw component
query still exhausts the shared allowance. The 32-crossing stabilized
circle also remains capped during its raw query. These observations locate
remaining costs; they do not turn zero nonunit pivots in this corpus into
a theorem for arbitrary triangulations.

\subsection{Serial measurements of complete operations}

Both experiments run serially after the tests, source audit and runtime
publication. Each uses five shuffled paired rounds, old and new A/A arms,
and an excluded warmup per arm and input. All 220 measured calls and
44 warmups finish. Tables report individual-arm medians and medians of
within-round ratios; the latter need not equal ratios of the former.

The preparation endpoint starts from a fresh PD, constructs the canonical
exterior, independently replays its source geometry, obtains the primitive
cocycle and validates all normal coordinates on the reconstructed manifold.
Only the cocycle solver changes. This includes construction and validation
costs rather than timing an isolated elimination loop.

\begin{center}
\small
\begin{tabular}{lrrrrr}
\toprule
Input & Old (ms) & New (ms) & Old/new & Old A/A & New A/A \\
\midrule
circle-1 & 4.901 & 5.005 & 0.969 & 0.998 & 1.006 \\
circle-8 & 56.628 & 41.304 & 1.371 & 1.000 & 1.030 \\
circle-32 & 247.496 & 186.966 & 1.343 & 0.989 & 1.073 \\
circle-64 & 508.502 & 368.708 & 1.372 & 1.012 & 0.996 \\
Survivor 00 & 113.154 & 81.171 & 1.400 & 1.002 & 0.999 \\
Gordian & 2244.886 & 1014.818 & 2.222 & 0.992 & 0.998 \\
\bottomrule
\end{tabular}

\end{center}

All 120 measured calls and 24 warmups finish in 59.152 seconds of driver
elapsed time. Every old and new repetition produces the same seed hash.
The larger circles and survivor input improve by $1.343$--$1.400$; Gordian
improves by $2.222$, with medians falling from 2.244886 to 1.014818 seconds.
Old A/A ratios range from $0.989$ to $1.012$, and new A/A from $0.996$ to
$1.073$. The one-crossing input has paired ratio $0.969$: about three
percent overhead, with individual medians 4.901 and 5.005 milliseconds.
Maintaining all row buckets and incidence sets has a small-input cost.

The recognition endpoint uses the same current recognizer in both arms
and swaps only primitive cocycle discovery. Earlier shortcuts are disabled
so that the native stage is reached. A raw seed proves both stabilized
circles, the five-crossing example needs span optimization, and the two
remaining cases continue to the exact Khovanov fallback. Every verdict
matches the known input. Positive native paths include their normal
component query and independent source-proof replay.

\begin{center}
\small
\begin{tabular}{lrrrrr}
\toprule
Input & Old (ms) & New (ms) & Old/new & Old A/A & New A/A \\
\midrule
circle-8 & 388.982 & 375.429 & 1.038 & 0.999 & 1.035 \\
circle-16 & 1446.584 & 1428.088 & 1.008 & 0.998 & 1.002 \\
Optimized positive & 295.827 & 302.754 & 1.042 & 0.957 & 0.930 \\
Genus-one miss & 170.545 & 168.546 & 0.990 & 1.007 & 1.008 \\
Trefoil & 90.172 & 94.203 & 1.022 & 1.000 & 1.004 \\
\bottomrule
\end{tabular}

\end{center}

All 100 measured recognitions and twenty warmups finish in 60.576 seconds.
Paired old/new ratios range from $0.990$ to $1.042$, old A/A from $0.957$
to $1.007$, and new A/A from $0.930$ to $1.035$. These results do not
establish a broad whole-recognition gain. In particular, the optimized
positive's apparent paired gain is smaller than its control fluctuations;
its separate arm medians even reverse order. Component queries, span
optimization and fallback dominate these complete calls. The native option
remains disabled by default, and this experiment makes no new claim about
the ordinary portfolio that decides its inputs before reaching this stage.


The reproducible driver is
\path{../fast/normal_orbit_research/cocycle_sparse.py}, with \code{audit},
\code{prepare} and \code{recognize} modes. Audit and measurements retain
470 matching source and fixture hashes, including the prior full source
certificate audit. The two archives still in the incoming directory are covered
by their earlier integration records. The checked \code{openai/math} head
remains \code{fd4aeeb2ee4fc729c18d98444fed42fd0529eeeb}; this change uses
exact sparse elimination, not a new external complexity claim.

